\section*{Chapter \ref{ch:m1:financing} --- Financement~: repo, prêt de titres et prime brokerage}

\begin{solution}{exo:m1:financing:1}
$L = 6$~; fonds propres anéantis par une baisse de $1/6 = 16.7\%$~; $r_E = 6 \times 6\% - 5
\times 4\% = 16\%$.
\end{solution}

\begin{solution}{exo:m1:financing:2}
Trésorerie levée $0.97 \times 50 = 48.5$ millions de dollars~; intérêts $48.5\times10^6
\times 0.042 \times 7/360 = \$39\,608$.
\end{solution}

\begin{solution}{exo:m1:financing:3}
Le \omterm{def:m1:financing:seclending}{taux de rebate} vaut $4\% - 6\% = -2\%$~: la \omterm{def:m1:financing:short}{vente à découvert} coûte 400\,000~dollars par an en carry.
\end{solution}

\begin{solution}{exo:m1:financing:4}
Prêt de 48~dollars. \omterm{def:m1:financing:margincall}{Appel de marge} à $48/(1-0.30) = \$68.57$. À 60~dollars, les fonds propres valent 12~dollars
contre un minimum exigé de $0.40\times60 = \$24$~: 12~dollars par action.
\end{solution}

\begin{solution}{exo:m1:financing:5}
Actifs de 2.85~milliards de dollars, fonds propres $0.5 - 0.15 = 0.35$ milliard de dollars~: $L = 8.14$.
Pour revenir à 6~: vendre $5 \times 0.05 \times 3 = 0.75$ milliard de dollars. Pour revenir à 4~: les actifs doivent
être de 1.4~milliards de dollars, donc vendre 1.45~milliards de dollars, soit la moitié du portefeuille.
\end{solution}

\begin{solution}{exo:m1:financing:6}
$1 - 2.5/4 = 37.5\%$ du portefeuille. Rendement des fonds propres~: $0.12\% \times 40 =
4.8\%$ avant, $0.12\% \times 25 = 3.0\%$ après --- et la vente forcée aura
fait bouger les spreads mêmes dont le fonds est long.
\end{solution}

\begin{solution}{exo:m1:financing:7}
La dichotomie donne environ 0.129 par unité vendue. Dans une spirale linéarisée, chaque
tour, la baisse vaut $\lambda(L-1)$ fois la précédente, où $\lambda$ désigne le
coefficient d'impact~: la série diverge lorsque $\lambda \ge 1/(L-1) = 0.143$. Le seuil numérique
est un peu plus bas, car il n'est pas nécessaire que la série diverge~: il suffit
que la \emph{somme} des baisses atteigne $1/L = 12.5\%$ en douze
tours en partant de 5~\%.
\end{solution}

\begin{solution}{exo:m1:financing:8}
(i) La marge statique sur un prix doublé vaut 3.75~\% de la valeur courante, et non 7.5~\%.
(ii) L'horizon pertinent n'est pas un jour, mais le temps nécessaire pour liquider une
position concentrée, plusieurs jours pendant lesquels les ventes propres du dealer
font bouger le prix. (iii) Le pire jour historique a été observé alors que le
client achetait~; le mouvement qui compte est celui qui survient lorsqu'il, et
tous les autres dealers qui détiennent la même couverture vendent~: la distribution passée
exclut précisément cet événement. (À noter~: cinq années de hausse d'une action disent peu
sur sa queue gauche.)
\end{solution}

\begin{solution}{pb:m1:financing:1}
\textbf{1.} $L = 28/4 = 7$~; fonds propres anéantis par une baisse de 14.3~\%.
\textbf{2.} $28/1.6 = \$17.5$ milliards~; marge de 1.75~milliards de dollars.
\textbf{3.} $1.75/28 = 6.25\%$.
\textbf{4.} 2.25~milliards de dollars de trésorerie libre détenue ailleurs. Aucun dealer ne peut voir la
position totale ni vérifier que cette trésorerie n'est grevée d'aucune charge~: chacun exige une marge sur
son cinquième, comme si c'était la totalité.
\textbf{5.} $28 \times 4.45\% = 1.246$ milliards de dollars.
\textbf{6.} $1.246 - 0.16 = 1.086$ milliards de dollars, soit 27.2~\% des fonds propres.
\textbf{7.} $1.086/28 = 3.88\%$.
\textbf{8.} Perte de 2.8~milliards de dollars~; fonds propres de 1.2~milliards de dollars~; levier $25.2/1.2 =
21$.
\textbf{9.} 2.8~milliards de dollars sont appelés. Le family office verse 2.25~milliards de dollars~:
450~millions de dollars par dealer contre 560~millions dus, laissant 110~millions
de dollars impayés chez chacun.
\textbf{10.} Perte $0.08 \times 25.2 = 2.016$ milliards de dollars~; fonds propres $-0.816$
milliard de dollars~: le family office est insolvable et le levier n'a plus aucun sens.
\textbf{11.} Valeur de la couverture $23.184/5 = \$4.64$ milliards. Perte due au titre des swaps
$560 + 403 = \$963$ millions~; détenus en contrepartie, $350 + 450 = \$800$ millions.
\textbf{12.} Déficit de 163~millions de dollars, plus $0.12 \times 4\,637 = \$556$
millions de perte de liquidation~: 720~millions de dollars.
\textbf{13.} $163 + 0.03 \times 4\,637 = 302$ millions de dollars.
\textbf{14.} 2.8~milliards de dollars au lieu de 1.75~milliards de dollars.
\textbf{15.} $T = 50$ jours~; impact $0.7\times0.03\times\sqrt5 = 4.7\%$~; risque
$0.03\sqrt{50/3} = 12.2\%$~; \omterm{def:m1:financing:repo}{décote} $4.7 + 1.645\times12.2 = 24.8\%$, quatre
fois les 6.25~\% détenus.
\textbf{16.} Quatre autres dealers liquident les mêmes actions au même
moment, si bien que le volume disponible pour chacun n'est qu'une fraction du total, et
l'impact s'additionne~; et la marche aléatoire de la formule n'a pas de dérive, alors qu'ici
tout le marché sait que 23~milliards de dollars sont à vendre.
\textbf{17.} La position \emph{totale} du client auprès des dealers. Chacun aurait
vu une position représentant vingt-cinq jours de volume, et non cinq.
\textbf{18.} Ses ventes ont été réalisées avant que celles des autres n'aient fait bouger le
prix~; le dernier à vendre a vendu dans l'impact de tous les autres. Dans une
liquidation commune, l'ordre de sortie est tout le résultat, ce qui explique que les accords informels
de vente ordonnée ne tiennent pas.
\textbf{19.} \textbf{14.3~\%.}
\textbf{20.} Le levier mesure combien le client perd par unité de
mouvement du prix~; la marge mesure la part de cette perte que le prêteur a reçue
à l'avance.
\end{solution}

\begin{solution}{iq:m1:financing:1}
Obtenir un \omterm{def:m1:financing:seclending}{locate} auprès du \omterm{def:m1:financing:pb}{prime broker}. Vendre les actions sur le marché. Au
règlement-livraison, le \omterm{def:m1:financing:pb}{prime broker} emprunte les actions auprès d'un prêteur et les livre
à l'acheteur~; le produit de la vente va au prêteur à titre de collatéral en espèces,
réévalué au prix du marché chaque jour. Tant que la position reste ouverte~: percevoir des intérêts sur le collatéral au
\omterm{def:m1:financing:seclending}{taux de rebate} (taux moins les \omterm{def:m1:financing:seclending}{frais d'emprunt}), verser au prêteur tout dividende. Pour clôturer~: acheter
les actions, les restituer au prêteur, récupérer le collatéral~; le gain est
la baisse du cours, augmentée du rebate et diminuée des dividendes.
\iqlookfor{que le vendeur ne détient jamais le produit de la vente, et le dividende.}
\end{solution}

\begin{solution}{iq:m1:financing:2}
Une vente de titres assortie d'un engagement de rachat ultérieur à un prix plus
élevé~: un prêt de liquidités collatéralisé. Il est sûr pour le prêteur car il
détient un titre liquide d'une valeur supérieure au prêt (la \omterm{def:m1:financing:repo}{décote}), réévalué au
prix du marché chaque jour, et qu'il peut, dans la plupart des juridictions, vendre
aussitôt si l'emprunteur fait défaut, sans attendre un tribunal de faillite.
\iqlookfor{\omterm{def:m1:financing:repo}{décote}, \omterm{def:m1:financing:margincall}{appels de marge} quotidiens et droit de liquidation.}
\end{solution}

\begin{solution}{iq:m1:financing:3}
Vente $= (L-1)\,x\,A = 3 \times 5\% = 15\%$ de ses actifs. Vérification~:
actifs $100 \to 95$, fonds propres $25 \to 20$, actifs cibles 80, vente de 15.
\iqlookfor{le raccourci $(L-1)x$ ou une vérification rapide sur des nombres ronds.}
\end{solution}

\begin{solution}{iq:m1:financing:4}
Spread de financement sur le débit long (environ un demi-point de pourcentage
sur la partie financée)~; perte d'intérêts et spread de crédit sur le produit
des \omterm{def:m1:financing:short}{ventes à découvert}~; \omterm{def:m1:financing:seclending}{frais d'emprunt}, négligeables pour les grandes valeurs
liquides et de plusieurs pour cent pour les \omterm{def:m1:financing:short}{ventes à découvert} très demandées,
qui sont souvent précisément celles que le modèle recherche~; dividendes versés
sur les positions vendeuses, si le backtest a utilisé des rendements de cours~;
et effets de la retenue à la source sur les dividendes des positions longues. Pour un livre $3{:}2$ sur une unité de capitaux propres, les premiers postes seuls représentent environ 2~\% des capitaux propres par an~; avec des \omterm{def:m1:financing:short}{ventes à découvert} difficiles à emprunter, plusieurs fois ce montant.
\iqlookfor{\omterm{def:m1:financing:seclending}{frais d'emprunt} corrélés au signal.}
\end{solution}

\begin{solution}{iq:m1:financing:5}
Pour le client~: aucune inscription au registre des actions et, dans certaines juridictions,
aucune déclaration des participations importantes~; aucun droit de timbre ni taxe sur les transactions là où ils s'appliquent~;
marge négociée~; accès à des marchés où il ne peut pas détenir les actions directement.
Pour la banque~: un swap est un dérivé à son bilan, avec un traitement différent en matière de fonds propres et de compensation~;
elle perçoit le spread sur la totalité du notionnel et contrôle la couverture~;
en contrepartie, elle supporte le risque de crédit du client avec la seule marge contractuelle,
et la position totale du client est plus difficile à suivre.
\iqlookfor{les motivations des deux parties et le risque de crédit que prend le dealer.}
\end{solution}

\begin{solution}{iq:m1:financing:6}
Sur le coût de liquidation, et non sur la volatilité quotidienne~: marge $\approx$ impact plus un quantile élevé
du mouvement de prix sur les jours nécessaires pour vendre la position à un taux de participation réaliste,
avec des suppléments pour la concentration par valeur et pour le manque de diversification du portefeuille.
La rendre dynamique, sur la valeur courante. Exiger la déclaration des positions globales ailleurs
comme condition du crédit, et plafonner la limite si elle est refusée. Soumettre l'ensemble du portefeuille
à un test de résistance pour une baisse conjointe des valeurs corrélées. Veiller à ce que le contrat
permette de relever la marge à bref délai --- et exercer ce droit avant la crise,
car pendant celle-ci le client
ne peut pas payer.
\iqlookfor{les jours de liquidation comme horizon, et la demande de la position
globale.}
\end{solution}
